\documentclass[11pt]{article}
\usepackage{fontspec}
\usepackage[french,english]{babel}
\usepackage{amsmath}
\usepackage{amssymb}
\usepackage{amsthm}
\usepackage{mathrsfs}
\usepackage{tikz-cd}
\usepackage{enumitem}

\newcommand{\printedpage}[1]{\marginpar{\footnotesize p.~#1}}
\newcommand{\msp}[1]{\marginpar{\footnotesize [#1]}}
\newcommand{\uHom}{\underline{\mathrm{Hom}}}
\newcommand{\Hom}{\mathrm{Hom}}
\newcommand{\Ind}{\mathrm{Ind}}
\newcommand{\Ob}{\operatorname{Ob}}
\newcommand{\Fl}{\operatorname{Fl}}
\newcommand{\card}{\operatorname{card}}

\setlength{\parindent}{0pt}

\begin{document}

\printedpage{19}
% head: A. Grothendieck : Dérivateurs, Ch. XVIII, Édition des Universités de Montpellier II et Paris VII
% foot: Version 17 décembre 2025

\msp{page 84}
\[
\begin{tikzcd}
\uHom(J,C) \arrow[r, "\alpha^J"] \arrow[d, "\varepsilon_j^*"] & \uHom(J,\Ind'_\pi C) \arrow[d, "{\varepsilon_j^{*\prime}}"] & \ni \varphi \arrow[d, mapsto] \\
C \arrow[r, "\alpha"] & \Ind'_\pi(C) & \ni \varphi(j)\ ,
\end{tikzcd}
\qquad (\text{\footnotemark[21]})
\tag*{(4.4.6.2)}
\]
\footnotetext[21]{The functors $\varepsilon_j^*$, $\varepsilon_j^{*\prime}$ are the evaluation functors at $j$.}
which induces a functor
\[
\underbrace{\uHom(J/C)/\varphi}_{\overset{\text{def}}{=}\, I} \xrightarrow{\varepsilon_j^\varphi} C/\varphi(j)\ . \tag*{(4.4.6.3)} % [sic] \uHom(J/C) -> \uHom(J,C), everywhere else written Hom(J,C)
\]
If $C$ is stable under $\varinjlim$ relative to a set $K$ of the type of conditions $\leqslant \pi$ (e.g. finite $\varinjlim$), the same holds for $\uHom(J,C)$, and also for $\uHom(J,C)/\varphi$, and for $C/\varphi(j)$, by virtue of the fact that $\alpha^J$ and $\alpha$ commute with the $\varinjlim$ of the type envisaged, moreover (4.4.6.3) commutes with the limits in question. Let us therefore suppose $C$ stable at least under finite $\varinjlim$, so the source and target categories in (4.4.6.3) are filtering. (The fact that the functor commutes with finite $\varinjlim$ does not seem useful in the present context.) This will serve us to have the convenient cofinality criterion SGA 4 I 8.1.3 b) in terms of F\,1, F\,2.

\msp{page 85}
Let us recall that the cofinality of the functors (4.4.6.3), when it holds ( ?), implies the fact that $\varphi$ is a direct limit of the composite functor
\[
\begin{tikzcd}[row sep=small]
I \arrow[r, "\psi_J"] \arrow[d, equal] & \uHom(J,C) \arrow[r, "\alpha^J"] \arrow[dd, dotted, "\varepsilon_j^*"] & \uHom(J,\Ind'_\pi(C)) \arrow[dd, dotted, "{\varepsilon_j^{*\prime}}"] \\
\uHom(J,C)/\varphi \arrow[d, "\varepsilon_j^\varphi"'] & & \\
C/\varphi(j) \arrow[r, dotted, "\psi_j"] & C \arrow[r, dotted, "\alpha"] & \Ind'_\pi(C)\ ,
\end{tikzcd}
\tag*{(4.4.6.4)}
\]
along the filtering ordered set $I$, more precisely that $\varphi$ is the $\varinjlim$ ``argument by argument in $j$'', which means that for every $j \in \Ob J$, $\varphi(j) = \varepsilon_j^*(\varphi)$ is a direct limit of $\varepsilon_j^*(\alpha^J\psi) = (\alpha\psi_j)\varepsilon_j^\varphi$. (Now we do indeed know that $\varphi(j)$ is the limit of $\alpha\psi$, since $\varphi(j) \in \Ob C^\wedge$ is ind-representable, hence also that of $(\alpha\psi_j)\varepsilon_j^\varphi$ if $\varepsilon_j^\varphi$ is cofinal.)

Thus, admitting the cofinality of (4.4.6.3) and interpreting $\psi_J$ in (4.4.6.4)

\msp{page 86}
\[
\psi_J \ :\ I \longrightarrow \uHom(J,C)
\]
as a functor
\[
I \times J \xrightarrow{\Psi} C\ , % [sic] \Psi -> \psi, the functor I x J -> C is called \psi in the next formula and in (4.4.7.1)
\]
we find the ``indicial representation'' of $\varphi$
\[
\varphi(j) \ =\ \underbrace{\varinjlim_{i\in I} \psi(i,j)}_{\substack{\text{direct limit}\\ \text{taken in } \Ind'_\pi(C)\\ \text{of objects of } C \hookrightarrow \Ind'_\pi(C)}} \qquad ,
\]

\printedpage{20}
% head: A. Grothendieck : Dérivateurs, Ch. XVIII, Édition des Universités de Montpellier II et Paris VII
% foot: Version 17 décembre 2025
whence
\[
\begin{aligned}
\underbrace{\varinjlim_{j\in J}{}_{\,\Ind'_\pi C}\varphi(j)}_{\text{subject to existence}} \quad &= \quad \varinjlim_{j\in J}\, \varinjlim_{i\in I} \psi(i,j)\\
&= \quad \varinjlim_{i\in I}\, \varinjlim_{j\in J} \psi(i,j)\ ,
\end{aligned}
\]
all direct limits being taken in $\Ind'_\pi(C)$. But if the $\varinjlim$ of type $J$ are representable in $C$, and $\card J \leqslant \pi$, we know that the $\varinjlim\limits_{j\in J} \psi(i,j)$ \textit{in} $C$ is also their $\varinjlim$ in $\Ind'_\pi C$, and we therefore find that $\varinjlim\limits_{j\in J} \varphi(j)$ exists in $\Ind'_\pi(C)$ if and only if $\varinjlim\limits_{i\in I} \Psi(i)$ exists in $\Ind'_\pi(C)$, where for $i \in I$, we set
\[
\Psi(i) \ =\ \varinjlim_{j\in J} \psi(i,j) \qquad \textit{in } C.
\]

\msp{page 87}
Let us summarize what we have obtained:

\textbf{Lemma 4.4.7.} \textit{Let $C$ be a small category, $\pi$ an infinite cardinal, $J$ a category such that $\card \Fl J \leqslant \pi$, $\varphi : J \longrightarrow \Ind'_\pi(C)$ a functor. Suppose the $\varinjlim$ of type $J$ to exist in $C$, and that the functors $\varepsilon_j^\varphi$ ($j \in \Ob J$) of (4.4.6.3) are \textup{cofinal} (which is made explicit in a simple way especially if source and target are filtering, something established if $C$ is stable under finite $\varinjlim$). Then $\varphi$ has a ``canonical indicial representation''}
\[
\psi \ :\ I \times J \longrightarrow C \tag*{(4.4.7.1)}
\]
\textit{with}
\[
I \ =\ \uHom(J,C)/\varphi\ , \tag*{(4.4.7.2)}
\]
\textit{i.e. we have the isomorphism functorial in $j \in \Ob J$}
\[
\varphi(j) \ \simeq\ \varinjlim_{I} \psi(i,j) \qquad (\text{\footnotemark[22]})\ . \tag*{(4.4.7.3)}
\]
\footnotetext[22]{direct limit in $\Ind'_\pi$ of objects of $C$.}
\textit{Let us set for every $i \in \Ob I$}
\[
\Psi(i) \ =\ \varinjlim_{J} \psi(i,j) \qquad (\text{\footnotemark[23]})\ . \tag*{(4.4.7.4)}
\]
\footnotetext[23]{direct limit taken in $C$, but valid also in $\Ind'_\pi C$.}
\textit{Then $\varinjlim\limits_{j\in J} \varphi(j)$ (in $\Ind'_\pi C$) exists if and only if}
\msp{page 88}
\textit{the $\varinjlim\limits_{i\in I} \Psi(i)$ exists (in $\Ind'_\pi C$ as well), and the two are equal.}

\textbf{Corollary 4.4.8.} \textit{Suppose that $C$ is stable under $\varinjlim$ of cardinality $\leqslant \pi$ (not only under those of type $J$). Then $\uHom(J,C)$ and $I = \uHom(J,C)/\varphi$ have the same property, $I$ is large relative to $\pi$, and $\varinjlim\limits_{j\in J} \varphi(j) \ =\ \varinjlim\limits_{i\in I} \Psi(i)$ exists in $\Ind'_\pi(C)$, subject to cofinality of the functors (4.4.6.3).}

\printedpage{21}
% head: A. Grothendieck : Dérivateurs, Ch. XVIII, Édition des Universités de Montpellier II et Paris VII
% foot: Version 17 décembre 2025
Indeed it suffices to verify this

\textbf{Lemma 4.4.8.1.} \textit{Let $I$ be a category large relative to $\pi$, and $\Psi : I \longrightarrow C$ a functor. Then $\varinjlim\limits_{i} \Psi(i)$ exists in $\Ind'_\pi(C)$, and it is the direct limit in $C^\wedge$, or again ``it is'' the ind-object large relative to $\pi$ represented by $\Psi$.} (In other words (\footnote[24]{This is an abuse, this stability is a distinctly stronger statement! I prove it on page 99.}), $\Ind'_\pi(C) \subset C^\wedge$ is a strictly full subcategory \textit{stable} under $\varinjlim$ large relative to $\pi$.)

\msp{page 89}
This lemma is a tautology: if $F$ denotes the direct limit in $C^\wedge$, it is clear by definition (taking a cofinal functor $I_0 \longrightarrow I$, with $I_0$ \textit{ordered} and large relative to $\pi$, so that the $\varinjlim$ can be taken along $I_0$) that $F$ is in $\Ind'_\pi(C)$, and then it is a fortiori a $\varinjlim$ in $\Ind'_\pi C$.

To verify the existence of the $\varinjlim$ of cardinality $\leqslant \pi$ in $\Ind'_\pi C$, in the case where $C$ itself is stable under the said limits, it therefore suffices, thanks to Corollary 4.4.8, to prove the cofinality of the functors (4.4.6.3). This will follow from the

\textbf{Lemma 4.4.9.} \textit{Let $C$ be a small category, $\pi$ an infinite cardinal, $J$ a small category such that for all $j, j' \in \Ob J$, we have $\card \Hom(j,j') \leqslant \pi$. Under these conditions, for every $j \in \Ob J$, the functor (4.4.6.3) \textup{admits a left adjoint,} hence a fortiori it is cofinal (since the categories colocalized over $C/\varphi(j)$ have an initial object, a fortiori they are $0$-connected).}

\msp{page 90}
\textsc{Proof of} 4.4.8. Let us return to the diagram (4.4.6.2), from which we deduce (4.4.6.3), we are going to rewrite it in a simplified way:
\[
\begin{tikzcd}
B \arrow[r, "\beta"] \arrow[d, shift left, "\varepsilon"] & B' \arrow[d, shift left, "\varepsilon'"] & \ni y'_0 \arrow[d, mapsto] \\
C \arrow[r, "\gamma"] \arrow[u, dotted, shift left, "\lambda"] & C' \arrow[u, dotted, shift left, "\lambda'"] & \ni z'_0\ .
\end{tikzcd}
\tag*{(4.4.9.1)}
\]
I say that $\varepsilon$, $\varepsilon'$ have \textit{left adjoints} $\lambda$ resp. $\lambda'$, and that the canonical ``base change'' homomorphism
\[
\lambda'\gamma \longrightarrow \beta\lambda \tag*{(4.4.9.2)}
\]
is an \textit{isomorphism.} Indeed, the existence of the adjoints of $\varepsilon$, $\varepsilon'$, namely of the functors $\varepsilon_{j\,!}$ and $\varepsilon'_{j\,!}$ relative to the inclusion
\[
\varepsilon_j \ :\ e \longrightarrow J
\]
and to the functor categories $\uHom(J,C)$ resp. $\uHom(J,\Ind'_\pi C)$, resulting from the existence of sums of type $\Hom(j,j')$ ($j' \in \Ob J$) in $C$ (already established), \textit{and in} $\Ind'_\pi C$\textit{ ;} and the fact that the functor $C \longrightarrow \Ind'_\pi C$ commutes
\msp{page 91}
with the said sums (something also already established, since $\card\, (\Hom(j,j')) \leqslant \pi$) then implies that (4.4.4.2) is an isomorphism. Subject to the pending verification (the existence % [sic] (4.4.4.2) -> (4.4.9.2), the isomorphism in question is the change-of-base arrow (4.4.9.2)
\printedpage{22}
% head: A. Grothendieck : Dérivateurs, Ch. XVIII, Édition des Universités de Montpellier II et Paris VII
% foot: Version 17 décembre 2025
of sums of cardinality $\leqslant \pi$ in $\Ind'_\pi(C)$), Lemma 4.4.9 will follow from the more general result:

\textbf{Lemma 4.4.10.} \textit{Consider a commutative square (4.4.9.1) (solid lines) of categories, where $\varepsilon$ admits a left adjoint $\lambda$, $\varepsilon'$ a left adjoint $\lambda'$, and the base change arrow (4.4.9.2) is an isomorphism. Then for every $y'_0$ in $B'$, setting $\varepsilon'(y'_0) = z'_0$, the functor induced by $(\varepsilon, \varepsilon')$,}
\[
\overline{B} \ =\ B/y'_0 \xrightarrow{\bar{\varepsilon}} \overline{C} = C/z'_0\ ,
\]
\textit{admits a left adjoint $\overline{\lambda}$, given by}
\[
\begin{tikzcd}[column sep=small]
\overline{\lambda}(\gamma(z) \xrightarrow{v} z'_0) \ \ = & \Bigl(\lambda'\gamma(z) \simeq \beta(\lambda(z)) \arrow[r, "{\lambda'(v)}"] \arrow[ddr, "{\substack{\overline{\lambda}(v) \,=\, \text{composite}\\ \text{adj} \circ \lambda'(v)}}"'] & \lambda'(z'_0) \arrow[d, equal] \\
 & & \lambda'\varepsilon'(y'_0) \arrow[d, "\text{adj.}"] \\
 & & y'_0\Bigr)\ \ .
\end{tikzcd}
\]

\msp{page 92}
The proof is carried out by a tedious AQT \texttt{[âne qui trotte: at a donkey's trot, i.e. by routine verification]} (I confess not to have pushed to the last detail \ldots).

To finish the proof of 4.4.9, there remains to prove one last ( ?)

\textbf{Lemma 4.4.11.} \textit{Let $C$ be a small category, stable under sums of cardinality $\leqslant \pi$. Then $\Ind'_\pi C$ is stable under sums of cardinality $\leqslant \pi$ (and, for the record, the functor $\alpha : C \longrightarrow \Ind'_\pi C$ commutes with them).}

Applying Corollary 4.4.8 to Lemma 4.4.7, we are reduced to proving the cofinality of (4.4.6.3) when $J$ is discrete, plus the fact that $\uHom(J,C)/\varphi$ is large relative to $\pi$, in this case. But giving $\varphi$ amounts to giving a family $(\varphi_j)_{j\in J}$ of objects of $\Ind'_\pi(C)$, and the functor (4.4.6.3) is made explicit as the projection
\[
\prod_{j\in J} (C/\varphi_j) \longrightarrow C/\varphi_j
\]
from the total product onto one of the factors. So we are reduced to this:

\msp{page 93}
\textbf{Lemma 4.4.12.} \textit{Let $(C_j)_{j\in J}$ be a family of small categories large relative to $\pi$, indexed by a small set $J$. Then $\prod C_j$ is large relative to $\pi$, and the projections (for $j_0 \in J$) $\prod C_j \longrightarrow C_{j_0}$ are cofinal.}

Let us first prove cofinality, which is the most important for us. This is obvious if $J = \{j_0\}$, otherwise we have $J = J' \cup \{j_0\}$, and $\prod_{j\in J} C_j \simeq C' \times C_{j_0}$, where $C' = \prod_{j\in J'} C_j$, so we must see that $C' \times C_{j_0} \longrightarrow C_{j_0}$ is cofinal, i.e. that every category $\xi\backslash(C' \times C_{j_0})$ (for $\xi \in \Ob C_{j_0}$) is 0-connected, now it is isomorphic to $C' \times (\xi\backslash C_{j_0})$, and $\xi\backslash C_{j_0}$ being 0-connected, it remains to see that $C'$ is 0-connected. We should be careful in this regard that an \textit{infinite} product of 0-connected categories is not necessarily 0-connected, so cofinality is not completely tautological! But if we prove that the category $\prod_{j\in J} C_j$ is
\printedpage{23}
% head: A. Grothendieck : Dérivateurs, Ch. XVIII, Édition des Universités de Montpellier II et Paris VII
% foot: Version 17 décembre 2025
\msp{page 94}
large relative to $\pi$, hence filtering, this implies that it is 0-connected, and we can apply this same result to $C' = \prod_{j\in J'} C_j$, and we win.

So we are reduced to proving that $\prod C_j$ is large relative to $\pi$. This can be seen either from the criterion $\mathrm{PF}_\pi\,1$, $\mathrm{PF}_\pi\,2$ of 4.4.4 b, or by reducing to the case where the $C_j$ are ordered, but that looks longer. Let us note that we prove in the same way that a (small) product of filtering categories is again a filtering category.

\textbf{Remarks.}
\begin{enumerate}[label={}, leftmargin=*]
\item[1.] This argument in no way shows (even if $C$ is stable under $\varinjlim$ of cardinality $\leqslant \pi$) that for every $\varphi \in \uHom(J,\Ind(C))$, the functor similar to (4.4.6.3)
\[
\uHom(J,C)/\varphi \longrightarrow C/\varphi(j)
\]
is cofinal, and I doubt that it always is. The difficulty comes from the fact that $C \longrightarrow \Ind\, C$ does not commute with sums of cardinality $\leqslant \pi$, only with finite sums,
\msp{page 95}
so Lemma 4.4.10 cannot be invoked.
\item[2.] We have just proved, via 4.4.8 and 4.4.9 (whose proof is completed with 4.4.12), that if $C$ is a small category stable under $\varinjlim$ of cardinality $\leqslant \pi$, then $\Ind_\pi C$ is also stable under these limits. Moreover, the direct systems of this type $J$ in $\Ind_\pi C$ admit an indicial representation
\[
\psi \ :\ I \times J \longrightarrow C\ ,
\]
with $I$ large relative to $\pi$. And we can always suppose, at the cost of taking $I' \longrightarrow I$ cofinal with $I'$ an ordered set large relative to $\pi$, that $I$ is moreover an ordered set.
\end{enumerate}

That is the means then, for the proof of the

\textbf{Theorem 4.4.13.} \textit{Let $C$ be a category equivalent to a small category, and stable under $\varinjlim$ of cardinality $\leqslant \pi$, where $\pi$ is a given infinite cardinal. Then:}
\msp{page 96}

\begin{enumerate}[label={}, leftmargin=*]
\item[\textit{a)}] \textit{(For the record.) The canonical functor}
\[
C \longrightarrow \Ind_\pi(C) \xrightarrow{\ \sim\ } \Ind'_\pi(C)
\]
\textit{commutes with $\varinjlim$ of cardinality $\leqslant \pi$.}
\item[\textit{b)}] \textit{(For the record.) The direct systems of cardinality $\leqslant \pi$ in $\Ind_\pi(C)$ (i.e. the functors $\varphi : J \longrightarrow \Ind_\pi(C)$, with $\card \Fl(J) \leqslant \pi$) admit an indicial representation}
\[
\psi \ :\ I \times J \longrightarrow C\ ,
\]
\textit{with $I$ an ordered set large relative to $\pi$, and admit a direct limit in $\Ind_\pi C$ which are computed by}
\[
\underbrace{\varinjlim_{j\in J} \varphi(j)}_{\text{in } \Ind_\pi(C)} \ \simeq\ \text{``}\,\varinjlim_{i\in I}\,\text{''}\Psi(i)
\]
\end{enumerate}

\end{document}
